\chapter{CONCLUSÃO}
\label{cap:conclusao}

A motivação para este estudo surgiu da necessidade crescente de desenvolver APIs robustas e eficientes que possam atender às exigentes demandas de desempenho e escalabilidade de sistemas de software atuais. Diante disso, 
foi conduzido um estudo empírico que emprega procedimentos estatísticos rigorosos para quantificar as melhorias oferecidas por uma linguagem de programação e um estilo arquitetural de comunicação considerados promissores dentro de um contexto de sistemas distribuídos, abordando especificamente Go e Java e gRPC e REST. 

A pesquisa foi impulsionada por duas principais questões de pesquisa: (i)~qual a combinação mais eficiente entre as linguagens de programação Go e Java e os estilos arquiteturais de comunicação gRPC e REST e (ii)~qual a influência da linguagem de programação e do estilo arquitetural de comunicação na eficiência. 

A questão de pesquisa~\#1 evidenciou que Go apresenta desempenho superior tanto em requisições \textit{StdSize} quanto \textit{LargeSize}, 
comparado ao Java. Esse resultado sugere que a seleção de Go em combinação com gRPC, pode oferecer vantagens substanciais em termos de eficiência, especialmente em ambientes que requerem comunicações frequentes e de alto desempenho.

A questão de pesquisa~\#2 conclui que -- de fato -- a influência da arquitetura de comunicação é de 94\% no desempenho.
\okterrasblp{
Isso indica que, ao substituir REST por gRPC, existem cenários em que fábricas de software podem ter uma melhoria de desempenho em seus sistemas 
sem necessitar de uma reescrita de código em uma linguagem de programação diferente.}

\noterrasblp{A diferença de desempenho entre REST e gRPC nos diferentes tamanhos de requisição pode ser atribuída a diversos fatores inerentes às duas arquiteturas de comunicação. O gRPC, por ser baseado no protocolo HTTP/2 e utilizar a serialização Protobuf, é mais eficiente em termos de latência e consumo de recursos para transmissões de dados menores devido à sua capacidade de multiplexação de múltiplas chamadas em uma única conexão. Já o REST, com sua simplicidade e ampla adoção, tende a ser mais eficiente para requisições maiores, possivelmente devido à sua capacidade de lidar com cargas de dados mais pesadas e menos dependência de conexões simultâneas.}

Trabalhos futuros incluem: 
(i)~avaliar a eficiência com requisições de diversos tamanhos a fim de encontrar pontos de transição onde uma arquitetura de comunicação pode superar outra; 
(ii)~avaliar a eficiência com diferentes frequências de solicitações, a fim de obter os limiares de desempenho de REST e gRPC em situações de estresse;
(iii)~avaliar em diferentes configurações de infraestrutura como rede, CPU, memória, etc.; e
\okterrasblp{(iv)~avaliar o desempenho com algoritmos específicos no servidor, a fim de medir o impacto de funcionalidades especificas das linguagens, como paralelismo e gerenciamento de memória.}
